\documentclass{article}
\usepackage{amsmath, amssymb}
\usepackage{booktabs}
\usepackage{graphicx}
\usepackage{hyperref}
\usepackage{xcolor}
\usepackage{listings}
\usepackage{algorithm}
\usepackage{algorithmic}

% Code listing settings
\lstset{
    language=Python,
    basicstyle=\ttfamily\small,
    keywordstyle=\color{blue},
    commentstyle=\color{gray},
    stringstyle=\color{orange},
    showstringspaces=false,
    frame=single,
    breaklines=true,
    backgroundcolor=\color{gray!10}
}

\title{CRISPYx NB-GLM IRLS Optimization:\\Pre-Converged Control Statistics and Numba JIT Fusion}
\author{CRISPYx Development Team}
\date{January 2026}

\begin{document}

\maketitle

\begin{abstract}
We present a two-phase optimization strategy for the Negative Binomial Generalized Linear Model (NB-GLM) in CRISPYx that achieves up to \textbf{170× speedup} on large-scale CRISPR screening datasets. Phase 1 eliminates redundant control cell computations by pre-converging the control intercept and caching sufficient statistics. Phase 2 fuses element-wise IRLS operations into parallel Numba kernels, achieving an additional 26--57× speedup on the inner loop. Combined, these optimizations reduce the projected runtime for the Replogle-GW-k562 dataset (9,866 perturbations, 2M cells) from approximately 27 days to under 2 hours.
\end{abstract}

\section{Problem Statement}

\subsection{The Computational Bottleneck}

In CRISPR perturbation screens, each perturbation is compared against a shared pool of control cells. For a dataset with:
\begin{itemize}
    \item $n_{\text{control}} = 75{,}328$ control cells
    \item $n_{\text{pert}} \approx 400$ perturbation cells (per perturbation)
    \item $G = 8{,}248$ genes
    \item $P = 9{,}866$ perturbations
\end{itemize}

The na\"ive NB-GLM implementation performs Iteratively Reweighted Least Squares (IRLS) separately for each perturbation. Each IRLS iteration requires computing:
\begin{align}
    \eta &= X\beta + \text{offset} \label{eq:eta}\\
    \mu &= \exp(\eta) \label{eq:mu}\\
    W &= \frac{\mu^2}{\mu + \alpha\mu^2} \label{eq:weights}\\
    z &= \eta + \frac{Y - \mu}{\mu} \label{eq:working_response}
\end{align}

For \emph{all} cells (control + perturbation). With 10 IRLS iterations per perturbation:
\[
\text{Total matrix operations} = 10 \times 9{,}866 \times 75{,}728 \times 8{,}248 \approx 6.2 \times 10^{13}
\]

\subsection{Root Cause: Redundant Control Computation}

The key insight is that \textbf{control cells are identical across all perturbations}, yet their contributions to $W$ and $Wz$ are recomputed for every perturbation. This redundancy dominates runtime because $n_{\text{control}} \gg n_{\text{pert}}$.

\section{Phase 1: Pre-Converged Control Sufficient Statistics}

\subsection{Mathematical Foundation}

For a two-coefficient model (intercept $\beta_0$ + perturbation effect $\beta_1$), the design matrix is:
\[
X = \begin{pmatrix}
1 & 0 \\
\vdots & \vdots \\
1 & 0 \\
1 & 1 \\
\vdots & \vdots \\
1 & 1
\end{pmatrix}
\quad
\begin{matrix}
\leftarrow \text{control cells} \\
\\
\\
\leftarrow \text{perturbation cells}\\
\\
\end{matrix}
\]

The normal equations for weighted least squares are:
\begin{align}
    X^\top W X \cdot \beta &= X^\top W z \\
    \begin{pmatrix}
        \sum W_i & \sum_{i \in \text{pert}} W_i \\
        \sum_{i \in \text{pert}} W_i & \sum_{i \in \text{pert}} W_i
    \end{pmatrix}
    \begin{pmatrix}
        \beta_0 \\ \beta_1
    \end{pmatrix}
    &=
    \begin{pmatrix}
        \sum W_i z_i \\
        \sum_{i \in \text{pert}} W_i z_i
    \end{pmatrix}
\end{align}

This decomposes into control and perturbation contributions:
\begin{align}
    X^\top W X &= \underbrace{\begin{pmatrix} W_{\text{ctrl}} & 0 \\ 0 & 0 \end{pmatrix}}_{\text{control}} + \underbrace{\begin{pmatrix} W_{\text{pert}} & W_{\text{pert}} \\ W_{\text{pert}} & W_{\text{pert}} \end{pmatrix}}_{\text{perturbation}} \\
    X^\top W z &= \underbrace{\begin{pmatrix} (Wz)_{\text{ctrl}} \\ 0 \end{pmatrix}}_{\text{control}} + \underbrace{\begin{pmatrix} (Wz)_{\text{pert}} \\ (Wz)_{\text{pert}} \end{pmatrix}}_{\text{perturbation}}
\end{align}

where $W_{\text{ctrl}} = \sum_{i \in \text{control}} W_i$ and $(Wz)_{\text{ctrl}} = \sum_{i \in \text{control}} W_i z_i$.

\subsection{The Frozen Control Approximation}

\textbf{Key observation:} If we pre-converge $\beta_0$ using only control cells, then $W_{\text{ctrl}}$ and $(Wz)_{\text{ctrl}}$ become \emph{constants} that can be cached and reused across all perturbations.

\begin{enumerate}
    \item \textbf{Pre-convergence step} (one-time, per dataset):
    \begin{itemize}
        \item Fit intercept-only model on control cells until convergence
        \item Cache: $\beta_0^*$, $\alpha^*$ (dispersion), $W_{\text{ctrl}}^*$, $(Wz)_{\text{ctrl}}^*$
    \end{itemize}
    
    \item \textbf{Per-perturbation IRLS} (using cached control):
    \begin{itemize}
        \item Initialize $\beta_0 = \beta_0^*$, $\beta_1 = 0$
        \item Each iteration: only compute $W_{\text{pert}}$, $(Wz)_{\text{pert}}$ on perturbation cells
        \item Combine with cached $W_{\text{ctrl}}^*$, $(Wz)_{\text{ctrl}}^*$ to solve 2×2 system
    \end{itemize}
\end{enumerate}

\subsection{Complexity Analysis}

\begin{table}[h]
\centering
\caption{Computational complexity per perturbation}
\label{tab:complexity}
\begin{tabular}{lcc}
\toprule
\textbf{Approach} & \textbf{Cells per iteration} & \textbf{Operations/pert} \\
\midrule
Na\"ive & $n_{\text{ctrl}} + n_{\text{pert}}$ & $75{,}728 \times 8{,}248 \times 10 = 6.2 \times 10^9$ \\
Frozen Control & $n_{\text{pert}}$ only & $400 \times 8{,}248 \times 10 = 3.3 \times 10^7$ \\
\midrule
\textbf{Speedup} & & $\mathbf{188\times}$ (theoretical) \\
\bottomrule
\end{tabular}
\end{table}

\subsection{Implementation}

The optimization is implemented in \texttt{fit\_batch\_with\_control\_cache()} via the \texttt{use\_frozen\_control=True} parameter:

\begin{lstlisting}
# Cache control sufficient statistics (one-time)
control_cache = ControlStatisticsCache(
    beta_intercept=beta_0_converged,
    control_xtwx_intercept=W_ctrl_sum,    # sum(W) for control
    control_xtwz_intercept=Wz_ctrl_sum,   # sum(W*z) for control
    control_dispersion=alpha_converged,
    ...
)

# Per-perturbation IRLS (fast path)
for iteration in range(max_iter):
    # Only compute on perturbation cells!
    irls_iteration_with_reduction(
        Y_pert, beta_0, beta_1, offset_pert, alpha,
        eta_out, mu_out, W_pert_sum, Wz_pert_sum
    )
    
    # Combine with cached control
    xtwx_00 = W_ctrl_sum + W_pert_sum
    xtwz_0 = Wz_ctrl_sum + Wz_pert_sum
    
    # Solve 2x2 system (Cramer's rule)
    ...
\end{lstlisting}

\section{Phase 2: Numba JIT-Compiled IRLS Kernels}

\subsection{Motivation}

Even after Phase 1, the IRLS iteration on perturbation cells involves multiple temporary arrays:
\begin{lstlisting}
# NumPy implementation (multiple memory allocations)
eta = beta_0 + beta_1 + offset          # temp array 1
mu = np.exp(eta)                        # temp array 2
var = mu + alpha * mu * mu              # temp array 3
W = mu * mu / var                       # temp array 4
z = eta + (Y - mu) / mu                 # temp array 5
W_sum = np.sum(W, axis=0)               # reduction
Wz_sum = np.sum(W * z, axis=0)          # temp + reduction
\end{lstlisting}

Each temporary array causes:
\begin{itemize}
    \item Memory allocation overhead
    \item Cache eviction (poor memory locality)
    \item Multiple passes over the data
\end{itemize}

\subsection{Fused Numba Kernel}

We fuse all operations into a single pass with inline reductions:

\begin{lstlisting}
@nb.njit(parallel=True, cache=True)
def irls_iteration_with_reduction(
    Y, beta_intercept, beta_pert_coef, offset, alpha,
    log_min_mu, min_mu,
    eta_out, mu_out, W_sum_out, Wz_sum_out
):
    n_cells, n_genes = Y.shape
    
    for j in nb.prange(n_genes):  # Parallel over genes
        w_sum = 0.0
        wz_sum = 0.0
        
        for i in range(n_cells):
            # Fused computation (no temp arrays!)
            eta = beta_intercept[j] + beta_pert_coef[j] + offset[i, j]
            eta = max(log_min_mu, min(eta, 20.0))
            mu = max(exp(eta), min_mu)
            var = mu + alpha[j] * mu * mu
            W = (mu * mu) / max(var, min_mu)
            z = eta + (Y[i, j] - mu) / max(mu, min_mu)
            
            # Store and accumulate
            eta_out[i, j] = eta
            mu_out[i, j] = mu
            w_sum += W
            wz_sum += W * z
        
        W_sum_out[j] = w_sum
        Wz_sum_out[j] = wz_sum
\end{lstlisting}

\subsection{Micro-Benchmark Results}

\begin{table}[h]
\centering
\caption{Phase 2 kernel micro-benchmarks (8,000 genes)}
\label{tab:phase2}
\begin{tabular}{cccc}
\toprule
\textbf{Cells} & \textbf{NumPy (ms)} & \textbf{Numba (ms)} & \textbf{Speedup} \\
\midrule
300 & 38.4 & 0.7 & $\mathbf{57\times}$ \\
500 & 76.4 & 1.7 & $\mathbf{46\times}$ \\
1000 & 193.4 & 7.4 & $\mathbf{26\times}$ \\
\bottomrule
\end{tabular}
\end{table}

The speedup is higher for smaller cell counts because:
\begin{itemize}
    \item Memory allocation overhead is a larger fraction of total time
    \item Numba eliminates all temporary allocations
    \item Better CPU cache utilization with single-pass algorithm
\end{itemize}

\section{Combined Results}

\begin{table}[h]
\centering
\caption{End-to-end performance improvement}
\label{tab:combined}
\begin{tabular}{lccc}
\toprule
\textbf{Dataset} & \textbf{Original} & \textbf{Phase 1} & \textbf{Phase 1+2} \\
\midrule
Adamson (7.6K ctrl) & 109 s/pert & 16.6 s/pert & $\sim$3.5 s/pert \\
Replogle-GW (75K ctrl) & 237 s/pert & 1.4 s/pert & $\sim$0.5 s/pert \\
\midrule
\multicolumn{4}{c}{\textbf{Replogle-GW-k562 Full Run (9,866 perturbations)}} \\
\midrule
Total time & $\sim$27 days & $\sim$4 hours & $\sim$\textbf{2 hours} \\
\bottomrule
\end{tabular}
\end{table}

\section{Limitations and Future Work}

\subsection{Current Limitations}

\begin{enumerate}
    \item \textbf{Frozen control approximation:} The intercept $\beta_0$ is not updated during perturbation fitting. This is accurate when perturbation cells are a small fraction of total cells, which is typically true in CRISPR screens.
    
    \item \textbf{Offset centering:} The current implementation still uses NumPy for offset-centered working response computation. A fully fused kernel would provide additional speedup.
    
    \item \textbf{Legacy path available:} Set \texttt{use\_frozen\_control=False} to use the original IRLS algorithm for validation or edge cases.
\end{enumerate}

\subsection{Phase 3: Gene Chunking (Not Implemented)}

For extremely memory-constrained environments, gene chunking could reduce peak memory by processing genes in batches. This was deemed low priority since Phase 1+2 already achieve practical runtimes.

\section{Conclusion}

The two-phase optimization strategy transforms NB-GLM from computationally prohibitive to practical for large-scale CRISPR screens:

\begin{itemize}
    \item \textbf{Phase 1} exploits the structure of perturbation experiments (shared control) to eliminate $>99\%$ of redundant computation
    \item \textbf{Phase 2} applies low-level optimizations (loop fusion, parallel JIT) to the remaining computation
    \item Combined speedup: \textbf{170×} on Replogle-GW-k562, \textbf{30×} on Adamson
\end{itemize}

The implementation is available in CRISPYx v0.8.0 and enabled by default.

\end{document}
